\author{OTS}
\documentclass[12pt]{mybib}
\setheaderlogo{logo.png}

% Absatzformatierung
\setlength{\parindent}{0pt}
\setlength{\parskip}{0.6em}
\renewcommand{\arraystretch}{1.8}
\renewcommand{\DFA}{12mm}
\renewcommand{\DZA}{1.2}

\newcommand{\schreiblinien}[1]{% 
  \foreach \n in {1,...,#1}{% 
    \noindent\makebox[\linewidth]{\rule{\linewidth}{0.4pt}} }% 
}

\title{2. Petrus}
\author{Lothar Schmid}
\date{}

\begin{document}
\mytitlepage
\tableofcontents
 \begin{auslegung}{2. Petrus 3, 1-3} 
       \large\textbf{Schlachter}\normalsize     
        \switchcolumn
       \large\textbf{Kommentar}\normalsize       
        \switchcolumn*    
   
        \hh{1}\person[b, p]{Geliebte}, dies \verbN[b, p]{ist} \konj[b, p]{nun} schon der zweite Brief, den ich euch \verbN[b, p]{schreibe}, \konj[b, p]{um} durch Erinnerung eure lautere Gesinnung \verbN[b, p]{aufzuwecken},
        \switchcolumn   
        \hh{1}\\
         \schreiblinien{2}
        \switchcolumn*  
   
        \hh{2}\konj[b, p]{damit} ihr an die Worte \verbN[b, p]{gedenkt}, die von den heiligen \person[b, p]{Propheten} \verbN[b, p]{vorausgesagt worden sind} \konj[b, p]{und} dessen, was euch der Herr \konj[b, p]{und} Retter durch uns, die \person[b, p]{Apostel}, \verbN[b, p]{aufgetragen hat}.
        \switchcolumn
        \hh{2}\\
         \schreiblinien{4}
        \switchcolumn*  
   
        \hh{3}\konj[b, p]{Dabei} \verbN[b, p]{sollt} ihr vor allem das \verbN[b, p]{erkennen}, \konj[b, p]{dass} am Ende der Tage \person[b, p]{Spötter} \verbN[b, p]{kommen werden}, die nach ihren eigenen Lüsten \verbN[b, p]{wandeln}
        \switchcolumn
        \hh{3}\\
         \schreiblinien{3}
        \switchcolumn* 
   
        \hh{4}\konj[b, p]{und} \verbN[b, p]{sagen}: \enquote{\konj[b, p]{Wo} \verbN[b, p]{ist} die Verheissung seiner Wiederkunft? \konj[b, p]{Denn} seitdem die \person[b, p]{Väter} \verbN[b, p]{entschlafen sind}, \verbN[b, p]{bleibt} \konj[b, p]{alles} so wie es von Anfang der Schöpfung \verbN[b, p]{gewesen ist}!}
        \switchcolumn
        \hh{4}\\
         \schreiblinien{4}
        \switchcolumn*
   
        \hh{5}\konj[b, p]{Dabei} \verbN[b, p]{übersehen} sie \konj[b, p]{aber} absichtlich, \konj[b, p]{dass} es schon vorzeiten Himmel gab \konj[b, p]{und} eine Erde aus dem Wasser heraus [entstanden] \konj[b, p]{und} inmitten der Wasser \verbN[b, p]{bestanden hat} \konj[b, p]{durch} das Wort Gottes;
        \switchcolumn
        \hh{5}\\
         \schreiblinien{4}
        \switchcolumn*
   
        \hh{6}\konj[b, p]{und} \konj[b, p]{dass} durch diese [Wasser] die damalige Erde \konj[b, p]{infolge} einer Wasserflut zugrunde \verbN[b, p]{ging}.
        \switchcolumn
        \hh{6}\\
         \schreiblinien{2}
        \switchcolumn*
   
        \hh{7}Die jetzigen Himmel \konj[b, p]{aber und} die Erde \verbN[b, p]{werden} \konj[b, p]{durch} \konj[b, p]{dasselbe} Wort \verbN[b, p]{aufgespart} \konj[b, p]{und} für das Feuer \verbN[b, p]{bewahrt} \konj[b, p]{bis} zum Tag des Gerichts \konj[b, p]{und} des Verderbens der gottlosen \person[b, p]{Menschen}.
        \switchcolumn
        \hh{7}\\
         \schreiblinien{4}
        \switchcolumn*
   
        \hh{8}Dieses eine \konj[b, p]{aber} \verbN[b, p]{sollt} ihr nicht \verbN[b, p]{übersehen}, \person[b, p]{Geliebte}, \konj[b, p]{dass} ein Tag bei dem Herrn \verbN[b, p]{ist} wie tausend Jahre, \konj[b, p]{und} tausend Jahre wie ein Tag!
        \switchcolumn
        \hh{8}\\
         \schreiblinien{4}
        \switchcolumn*
        \hh{9}Der Herr \verbN[b, p]{zögert} nicht die Verheissung \verbN[b, p]{hinaus}, \konj[b, p]{wie} etliche es für ein Hinauszögern \verbN[b, p]{halten}, \konj[b, p]{sondern} er \verbN[b, p]{ist} langmütig gegen uns, \konj[b, p]{weil} er nicht \verbN[b, p]{will}, \konj[b, p]{dass} jemand verloren \verbN[b, p]{gehe}, \konj[b, p]{sondern dass} jedermann Raum zur Busse \verbN[b, p]{habe}.
        \switchcolumn
        \hh{9}\\
         \schreiblinien{4}
        \switchcolumn*
   
        \hh{10}Es \verbN[b, p]{wird} \konj[b, p]{aber} der Tag des Herrn \verbN[b, p]{kommen} wie ein Dieb in der Nacht; \konj[b, p]{dann} \verbN[b, p]{werden} die Himmel mit Krachen \verbN[b, p]{vergehen}, die Elemente \konj[b, p]{aber} vor Hitze sich \verbN[b, p]{auflösen} \konj[b, p]{und} die Erde \konj[b, p]{und} die Werke darauf \verbN[b, p]{verbrennen}.
        \switchcolumn
        \hh{10}\\
         \schreiblinien{4}
        \switchcolumn*
   
        \hh{11} \konj[b, p]{Da} nun dies alles \verbN[b, p]{aufgelöst wird}, wie sehr \verbN[b, p]{solltet} ihr euch \verbN[b, p]{auszeichnen} \konj[b, p]{durch} heiligen Wandel \konj[b, p]{und} Gottesfurcht,
        \switchcolumn
        \hh{11}\\
         \schreiblinien{3}
        \switchcolumn*
   
        \hh{12}\konj[b, p]{indem} ihr das Kommen des Tages Gottes \verbN[b, p]{erwartet} \konj[b, p]{und} ihm \verbN[b, p]{entgegeneilt}, an welchem die Himmel sich in Glut \verbN[b, p]{auflösen} \konj[b, p]{und} die Elemente vor Hitze \verbN[b, p]{zerschmelzen werden}!
        \switchcolumn
        \hh{12}\\
         \schreiblinien{4}
        \switchcolumn*
   
        \hh{13}Wir \verbN[b, p]{erwarten} \konj[b, p]{aber} nach seiner Verheissung neue Himmel \konj[b, p]{und} eine neue Erde, in denen Gerechtigkeit \verbN[b, p]{wohnt}.
        \switchcolumn
        \hh{13}\\
         \schreiblinien{3}
        \switchcolumn*
   
        \hh{14}\konj[b, p]{Darum}, \person[b, p]{Geliebte}, weil ihr dies \verbN[b, p]{erwartet}, \konj[b, p]{so} \verbN[b, p]{seid} eifrig darum \verbN[b, p]{bemüht}, \konj[b, p]{dass} ihr als unbefleckt \konj[b, p]{und} tadellos vor ihm \verbN[b, p]{erfunden werden} in Frieden!
        \switchcolumn
        \hh{14}\\
         \schreiblinien{4}
        \switchcolumn*
   
        \hh{15}\konj[b, p]{Und} \verbN[b, p]{seht} die Langmut unseres Herrn als [eure] Rettung \verbN[b, p]{an}, wie auch unser geliebter Bruder \person[b, p]{Paulus} euch \verbN[b, p]{geschrieben hat} nach der ihm gegebenen Weisheit,
        \switchcolumn
        \hh{15}\\
         \schreiblinien{4}
        \switchcolumn*
   
        \hh{16}\konj[b, p]{so} wie auch in allen Briefen, \konj[b, p]{wo} er von diesen Dingen \verbN[b, p]{spricht}, \konj[b, p]{was} die \person[b, p]{Unwissenden} \konj[b, p]{und} \person[b, p]{Ungefestigten} \verbN[b, p]{verdrehen}, \konj[b, p]{wie} auch die übrigen Schriften, zu ihrem eigenen Verderben.
        \switchcolumn
        \hh{16}\\
         \schreiblinien{4}
        \switchcolumn*
   
        \hh{17}Ihr \konj[b, p]{aber}, \person[b, p]{Geliebte}, \konj[b, p]{da} ihr dies im Voraus \verbN[b, p]{wisst}, \konj[b, p]{so} \verbN[b, p]{hütet} euch, \konj[b, p]{dass} ihr nicht durch die Verführung der \person[b, p]{Frevler} \konj[b, p]{mit} \verbN[b, p]{fortgerissen werdet} \konj[b, p]{und} euren eigenen festen Stand \verbN[b, p]{verliert}!
        \switchcolumn
        \hh{17}\\
         \schreiblinien{3}
        \switchcolumn*
   
        \hh{18}\verbI[b, p]{Wachst} dagegen in der Gnade \konj[b, p]{und} in der Erkenntnis unseres Herrn \konj[b, p]{und} Retters Jesus Christus! Ihm sei die Ehre, \konj[b, p]{sowohl} jetzt als auch \konj[b, p]{bis} zum Tag der Ewigkeit! Amen
        \switchcolumn
        \hh{18}\\
         \schreiblinien{3}
    \end{auslegung}




\end{document}